%% sections/04-study.tex

\section{Study Design}
\label{sec:study}

\subsection{Corpus and Comparison Groups}

The full corpus is 7 campaigns, 49 sessions. We define three comparison groups:

\begin{description}
  \item[Resource campaigns (n=2):] Campaign 4 (Supply Die + HP carry-over) and Campaign 6
    (Cohesion Die + HP carry-over). Both have explicit degrading-resource mechanics.
  \item[Grief-adjacent baselines (n=3):] Campaigns 1--3 (grief-themed, covenant-restoration,
    faith/authority). These campaigns have HP carry-over (between-session, partial recovery)
    but no degrading-resource mechanic. They provide a baseline for campaigns with HP
    pressure but without explicit attrition tracking.
  \item[Light-pressure baselines (n=2):] Campaigns 5 and 7 (Thieves Guild, Party of Strangers).
    Both have HP carry-over (partial recovery between sessions) but no combat-intensive
    encounters creating sustained pressure. Supply Die equivalent is not implemented.
\end{description}

\subsection{Primary Outcome: Character Agency Scores}

The primary outcome variable is Character Agency (CA) score per session, 0--10 on the
rubric. CA scores higher scores when PCs make choices that matter --- when the session's
final state depended on specific PC choices the module anticipated but did not require.
Under rubric v1.9--v1.10 (applied to C4 and C6), the 10-band additionally recognizes
sustained behavioral pattern, independent-motivation convergence, and character-register
conversion as paths to CA 10.

CA score per session position is our measure of whether resource pressure amplifies or
suppresses character-driven engagement.

\subsection{Controlling for Rubric Version}

Campaign 4 and 6 were scored under rubric v1.9--v1.12 (which includes CA-specific
amendments that raise the ceiling for CA 9--10 scoring). Campaigns 1--3 were scored
under rubric v1.0--v1.4 (which had a narrower CA 10-band). To enable cross-campaign
CA comparison, we analyze the \textit{trajectory shape} (monotone increasing vs.\ stable
vs.\ late-rising) rather than absolute score levels.

\subsection{Confounds}

The primary confound is campaign archetype: resource campaigns (C4, C6) are also distinct
in their motivational structure (siege-resource; duty-spine) from grief-adjacent baselines.
We cannot cleanly separate resource mechanic effect from archetype effect. We address
this by comparing C4 and C6 to C5 and C7 (campaigns with similar non-grief archetypes
but no resource mechanics) as a secondary comparison.
